% problem-000067 1 配位化学与氧化态
\section*{1. 配位化学与氧化态}
\textit{下列为分问。}

\subsection*{1-1}
利⽤不同的催化剂，⼄烯催化氧化可选择性地⽣成不同产物。产物A可使⽯灰⽔变浑浊。写出与A相对分⼦质量相等的其他所有产物的分⼦式和结构简式。

\subsection*{1-2}
化合物 $\mathrm { C s A u C l _ { 3 } }$ 呈抗磁性。每个Cs+周围有12个Cl−离⼦，每个Cl−离⼦周围有5个⾦属离⼦。⾦离⼦处在Cl−围成的配位中⼼（也是对称中⼼）。写出该化合物中Au的氧化态及其对应的配位⼏何构型。

\subsection*{1-3}
2019年1⽉，嫦娥四号成功在⽉球着陆。探测器上的五星红旗由⼀类特殊的聚酰亚胺制成。以下为某种芳⾹族聚酰亚胺的合成路线：

（图：）

画出A、B、C的结构简式。
